\SetPageTheme{theory}
\PageHeader{Sinteza 1}{Rezumatul mecanismelor invatate}{S1 • P1}

\begin{checkpointbox}[Rezumat functional]
  Pana aici am construit un lant logic coerent:
  \begin{itemize}
    \item \texttt{/} da catul;
    \item \texttt{\%} da restul;
    \item \texttt{\% 10} izoleaza ultima cifra;
    \item \texttt{/ 10} elimina ultima cifra;
    \item \texttt{cat timp (n != 0)} parcurge toate cifrele;
    \item \texttt{ogl = ogl * 10 + uc} construieste oglinditul.
  \end{itemize}
\end{checkpointbox}

\begin{guidebox}[Autoevaluare]
  Bifeaza mental ce poti explica fara ajutor:
  \begin{itemize}
    \item de ce avem nevoie de \texttt{copie};
    \item ce face \texttt{uc};
    \item de ce bucla se opreste;
    \item de ce formula pentru \texttt{ogl} functioneaza.
  \end{itemize}
\end{guidebox}

\newpage
\SetPageTheme{theory}
\PageHeader{Sinteza 2}{Cum se leaga de ideea mare}{S1 • P2}

\begin{missionbox}[Explicatie pedagogica]
  Aceasta editie nu este importanta fiindca palindroamele apar des in viata reala. Ea este importanta fiindca invata o forma de citire algoritmica a informatiei.

  Exact acelasi tip de disciplina apare atunci cand un calculator verifica un cod, citeste un mesaj, valideaza o structura sau pregateste date pentru un sistem mai complex.
\end{missionbox}

\begin{solobox}[Reflectie]
  Scrie un paragraf scurt despre legatura dintre:
  \begin{itemize}
    \item citirea atenta a cifrelor;
    \item verificarea algoritmica;
    \item protectia si controlul informatiei.
  \end{itemize}
  \AnswerSpace{5}
\end{solobox}
